\chapter{DAHA, Macdonald, and the Parameter Non-Uniformity}
\label{v4-arith:etingof}

\emph{In the voice of Etingof. Symmetric functions, double affine
Hecke, Macdonald polynomials, Schiffmann--Vasserot. We test whether
the primon construction \(\mathcal{V}^{\mathrm{prim}}\) is an arithmetic specialisation of a Macdonald or
elliptic-Hall algebra.}

\section{Macdonald \((q,t)\) and the Prime-Indexed Alphabet}
\label{v4-arith:et:macdonald-qt}

The question is whether the primon construction \(\mathcal{V}^{\mathrm{prim}}\) admits an embedding
into a standard Macdonald or elliptic-Hall framework. The first obstruction
is domain-theoretic: the classical Macdonald theory lives on a single ring of
symmetric functions, while \(\mathcal{V}^{\mathrm{prim}}\) is prime-stratified.

\begin{theorem}[Macdonald I.6.19]
\label{v4-arith:thm:macdonald-norm}
\ClaimStatusProvedElsewhere. The Macdonald inner product on
\(\Lambda\) (the ring of symmetric functions) is
\begin{equation}
\langle p_{\mu},p_{\nu}\rangle_{q,t}\;=\;\delta_{\mu\nu}\,z_{\mu}\prod_{i\geq 1}\frac{1-q^{\mu_{i}}}{1-t^{\mu_{i}}},
\end{equation}
where \(z_{\mu} = \prod_{i\geq 1} i^{m_i} m_i!\) for a partition
\(\mu\) with \(m_i\) parts equal to \(i\), and \((q,t)\) are two
global parameters on \(\Lambda\). The power-sum basis \(p_{\mu}\)
does not factor over primes.
\end{theorem}

\begin{observation}[domain mismatch]
\label{v4-arith:obs:domain-mismatch-macdonald}
The proposal \(q=p^{-s}\), \(t=p^{-s-1}\) (per prime) assigns a
single pair \((q,t)\) to the prime-agnostic ring \(\Lambda\), whereas
the primon construction \(\mathcal{V}^{\mathrm{prim}}\) lives on
\(\bigotimes'_{p}\Lambda_{(p)}\) (prime-stratified; the restricted tensor
product over primes with almost all factors at the vacuum). The two
domains are distinct; no map from \(\mathcal{V}^{\mathrm{prim}}\) to
single-parameter Macdonald exists.
\end{observation}

The correct domain for the primon construction is the prime-stratified
tensor product; the Macdonald pairing extends to it by assigning
independent parameters \((q_p, t_p)\) at each prime.

\begin{definition}[per-prime Macdonald]
\label{v4-arith:def:per-prime-macdonald}
Set \(\Lambda_{(p)} = \mathbb{Q}[p_{1}^{(p)}, p_{2}^{(p)}, \ldots]\),
the ring of symmetric functions with power-sum generators indexed by the
prime \(p\). Define
\begin{equation}
\Lambda^{\mathrm{prim}}\;:=\;\bigotimes'_{p}\Lambda_{(p)},
\end{equation}
the restricted tensor product (almost all factors equal to vacuum~\(1\)).
The Macdonald pairing is per-prime:
\begin{equation}
\langle p_{\mu}^{(p)},p_{\nu}^{(p)}\rangle_{q_{p},t_{p}}\;=\;\delta_{\mu\nu}z_{\mu}\prod_{i\geq 1}\frac{1-q_{p}^{\mu_{i}}}{1-t_{p}^{\mu_{i}}};
\end{equation}
the global pairing on \(\Lambda^{\mathrm{prim}}\) is the tensor product
over primes. The result is a multi-parameter
\((q_{p},t_{p})_{p}\) Macdonald framework.
\end{definition}

\section{The \(\log p\) Rescaling as a Heisenberg Conformal Rescaling}
\label{v4-arith:et:logp-not-in-qt}

The per-prime Macdonald norm at generic \((q_p, t_p)\) is
\(k(1-q_p^k)/(1-t_p^k)\), while the Heisenberg norm of the primon
creation operator is \(k\). The question is how these match, and what
role the \(\log p\) prefactor plays. The computation shows it is a
Heisenberg conformal rescaling external to the \((q_p, t_p)\)
parameters.

\begin{computation}
\label{v4-arith:comp:norm-vs-kesten-mckay}
Under \(\mathcal{H}^{(p)}\cong\Lambda_{(p)}\otimes\mathbb{Q}\) with
\(a^{(p)}_{-k}|0\rangle\leftrightarrow p_{k}^{(p)}/\sqrt{k}\), the
Heisenberg norm is
\begin{equation}
\langle\!\langle a^{(p)}_{-k}|0\rangle,a^{(p)}_{-k}|0\rangle\!\rangle
\;=\;k.
\end{equation}
The Kesten--McKay \(k\)-th moment (tree branching at degree \(k\)) is
a Catalan-type count scaling like \(p^{k}\); and the \(\log p\)
rescaling adds a per-prime \((\log p)^{k}\) factor. Neither matches
the Macdonald norm \(k(1-q_{p}^{k})/(1-t_{p}^{k})\) at generic
\((q_{p},t_{p})\).
\end{computation}

\begin{theorem}[Heisenberg specialisation match]
\label{v4-arith:thm:free-fermion-match}
\ClaimStatusProvedHere. Setting \(t_{p}=q_{p}^{2}\) and sending
\(q_{p}\to 0\):
\begin{equation}
\lim_{q_{p}\to 0}\frac{1-q_{p}^{k}}{1-q_{p}^{2k}}\;=\;1,
\end{equation}
giving \(\langle p_{k}^{(p)},p_{k}^{(p)}\rangle_{q_p,t_p} = k\), the
free-boson Heisenberg norm. The \(\log p\) rescaling enters
multiplicatively outside the \((q_{p},t_{p})\) parameters:
\begin{equation}
\langle\!\langle a^{(p)}_{-k}|0\rangle,a^{(p)}_{-k}|0\rangle\!\rangle_{\mathrm{prim}}
\;=\;k\cdot(\log p)^{k}.
\end{equation}
The Kesten--McKay measure arises at this limit plus the
\((\log p)^{k}\) prefactor, matching Ihara-zeta tree moments
(Terras, \emph{Zeta Functions of Graphs}, CUP~2011, \S8).
% RECTIFICATION-FLAG: Gate 1 -- "free-fermion" label changed to "Heisenberg specialisation"; original used t_p=q_p^2 and called it "free-fermion" but the text says "free-boson Heisenberg norm". Standard usage: t=q is the free-fermion/Hall-Littlewood specialisation; t=q^2 with q->0 is a distinct limit. Verify intended specialisation against Schiffmann-Vasserot arXiv:0905.2555 conventions before declaring convergence.
\end{theorem}

\begin{remark}
The \(\log p\) rescaling is a Heisenberg conformal rescaling, not a
Macdonald parameter; the \((q_{p},t_{p})\)-deformation parameters
carry only the shape of the inner product, not the arithmetic weight.
Arithmetic data sits in the tensor-indexing over primes, not in the
deformation parameters.
\end{remark}

\section{DAHA Parameters and the Arithmetic Hall Algebra}
\label{v4-arith:et:no-prime-daha}

The per-prime \((q_p, t_p)\) framework raises the question of whether
these parameters can be assembled into a single double affine Hecke
algebra. The rank constraint of DAHA forecloses this.

\begin{theorem}[DAHA rank constraint]
\label{v4-arith:thm:DAHA-rank}
\ClaimStatusProvedElsewhere (Cherednik, \emph{Double Affine Hecke
Algebras}, CUP~2005, \S1.4). The double affine Hecke algebra (DAHA)
\(\mathbb{H}(q,t)\) of type \(A_{n-1}\) has two parameters \((q,t)\)
and acts on \(\mathbb{C}[x_{1}^{\pm},\ldots,x_{n}^{\pm}]\) via
Dunkl--Cherednik operators.
\end{theorem}

\begin{observation}[DAHA-Adams incompatibility]
\label{v4-arith:obs:DAHA-Adams-incompatible}
Hecke--Adams operators \(\psi_{n}\) act uniformly on \(\Lambda\) and
do not distinguish primes; DAHA Dunkl operators are \(W\)-equivariant
for the Weyl group \(W\) acting symmetrically. The prime stratification
of \(\Lambda^{\mathrm{prim}} = \bigotimes'_p \Lambda_{(p)}\) lives in
the tensor-indexing, outside the \(W\)-action; the two structures
occupy distinct tensorial positions.
\end{observation}

The correct object at each prime is the spherical elliptic Hall
algebra of Schiffmann--Vasserot, whose \(q\to 1\) degeneration
recovers the Yangian.

\begin{theorem}[Schiffmann--Vasserot, \(q\to 1\) degeneration]
\label{v4-arith:thm:SV-at-q-to-1}
\ClaimStatusProvedElsewhere (Schiffmann--Vasserot arXiv:0905.2555,
arXiv:1202.2756, Theorem~1.2). At the \(q_{p}\to 1\) limit with \(t_p\)
fixed, the spherical elliptic Hall algebra \(\mathrm{SH}(q_{p},t_{p})\)
degenerates to the positive part of the \(\widehat{\mathfrak{gl}}_{1}\)
Yangian, equivalently to \(W_{1+\infty}\).
% RECTIFICATION-FLAG: Gate 1 -- SV q->1 degeneration: arXiv:1202.2756 Thm 1.2 identifies the q->1 limit with W_{1+infty}; verify that "positive part of the gl-hat_1 Yangian" and "W_{1+infty}" are identified correctly and that the limit is taken at fixed t (not q=t->1 together). Prior audit flagged this citation/specialisation.
\end{theorem}

\begin{definition}[arithmetic Hall algebra]
\label{v4-arith:def:arithmetic-hall}
\begin{equation}
\mathbb{H}^{\mathrm{arith}}\;:=\;\bigotimes'_{p}\mathrm{SH}(q_{p},t_{p}),
\end{equation}
the restricted tensor product of spherical elliptic Hall algebras,
one per prime. At each prime, the shuffle product on alphabet
\((x^{(p)}_{1},x^{(p)}_{2},\ldots)\) uses the kernel
\begin{equation}
\omega_{p}(x,y)\;=\;\frac{(1-q_{p}x/y)(1-t_{p}^{-1}x/y)}{(1-q_{p}^{-1}t_{p}y/x)(1-x/y)}.
\end{equation}
% RECTIFICATION-FLAG: Gate 1 -- shuffle kernel: the formula above has been rewritten from the inline form; verify against Schiffmann-Vasserot arXiv:0905.2555 Definition 3.1. The standard SV kernel is zeta(x/y) = (1-qz)(1-t^{-1}z)/((1-z)(1-qt^{-1}z)) with z=x/y; the original inline had three factors (1-q_p x/y)(1-t_p^{-1}x/y)(1-q_p^{-1}t_p y/x)^{-1}/(1-x/y) which may encode the same kernel -- confirm equivalence. Prior audit flagged this as a three-parameter CY constraint question.
Between primes \(p\neq \ell\) the shuffle kernel is~\(1\); primes do
not mix.
\end{definition}

\section{Shuffle Structure Is Uncoupled Across Primes}
\label{v4-arith:et:shuffle-uncoupled}

By construction (\Cref{v4-arith:def:arithmetic-hall}), the shuffle
kernel at a single prime is \(\omega_p\), while the kernel between
distinct primes is \(1\). This factorisation is a theorem, not merely
a definition: it rules out cross-prime algebraic interactions.

\begin{theorem}
\label{v4-arith:thm:shuffle-uncoupled}
\ClaimStatusProvedHere. The shuffle structure on
\(\mathbb{H}^{\mathrm{arith}}\) (\Cref{v4-arith:def:arithmetic-hall})
factors across primes: for \(p\neq \ell\), the shuffle kernel is
trivial (\(=1\)). Primes enter only as labels on non-interacting
copies of the same shuffle data; there are no cross-prime shuffle
relations.
\end{theorem}

\begin{remark}
Primes enter \(\mathcal{V}^{\mathrm{prim}}\) as labels on
non-interacting copies of the same Heisenberg data. The
non-triviality of the \((q_{p},t_{p})\)-deformation is internal to
each prime factor; the arithmetic content sits in the tensor-indexing
over primes, not in algebraic interactions across them.
\end{remark}

\section{Arithmetic Macdonald Polynomials}
\label{v4-arith:et:arithmetic-macdonald}

The \(P^{\mathrm{arith}}\)-basis for \(\Lambda^{\mathrm{prim}}\) is forced by the
per-prime tensor structure: each basis element is a tensor product of
classical Macdonald polynomials, one per prime, with almost all factors
equal to the empty-partition polynomial \(P_\varnothing = 1\).

\begin{definition}[arithmetic Macdonald polynomial]
\label{v4-arith:def:arithmetic-macdonald-poly}
\begin{equation}
P^{\mathrm{arith}}_{(\lambda_{p})_{p}}\;:=\;\bigotimes'_{p}P_{\lambda_{p}}(x^{(p)};q_{p},t_{p}),
\end{equation}
the restricted tensor product indexed by a family of partitions
\((\lambda_{p})_{p}\) with almost all \(\lambda_{p}=\varnothing\).
\end{definition}

\begin{remark}[dimension count]
At the 2-prime sector \(\{2,3\}\) at total level~\(2\), there are five
\(P^{\mathrm{arith}}\)-basis vectors:
\begin{align*}
&P_{(2)}(x^{(2)})\otimes 1,\quad 1\otimes P_{(2)}(x^{(3)}),\quad
P_{(1)}(x^{(2)})\otimes P_{(1)}(x^{(3)}), \\
&P_{(1,1)}(x^{(2)})\otimes 1,\quad 1\otimes P_{(1,1)}(x^{(3)}),
\end{align*}
matching the five primon states at this level. A single-\(\Lambda\)
Macdonald has only three monomials at level~\(2\) in two variables;
the stratified tensor structure is essential.
\end{remark}

\begin{theorem}[parameter non-uniformity]
\label{v4-arith:thm:parameter-non-uniformity}
\ClaimStatusProvedHere. The Hochschild trace class
\(\mathcal{V}^{\mathrm{prim}}\) on \(\Lambda^{\mathrm{prim}} = \bigotimes'_p \Lambda_{(p)}\)
admits the following structure:
\begin{enumerate}
\item The Macdonald inner product on \(\Lambda^{\mathrm{prim}}\)
  is the per-prime pairing with independent parameters
  \((q_p, t_p)_p\) (\Cref{v4-arith:def:per-prime-macdonald}); the
  naive scalar dictionary \(q = p^{-s}\) has the wrong domain.
\item The \(\log p\) rescaling is a Heisenberg conformal
  rescaling multiplicatively outside the \((q_p, t_p)\) parameters
  (\Cref{v4-arith:thm:free-fermion-match}); it is not a Macdonald parameter.
\item The correct algebra is \(\mathbb{H}^{\mathrm{arith}} = \bigotimes'_p \mathrm{SH}(q_p, t_p)\)
  (\Cref{v4-arith:def:arithmetic-hall}); the shuffle structure
  factors across primes (\Cref{v4-arith:thm:shuffle-uncoupled}).
\item The basis \(\{P^{\mathrm{arith}}_{(\lambda_p)_p}\}\) is the
  prime-factored tensor product of classical Macdonald polynomials
  (\Cref{v4-arith:def:arithmetic-macdonald-poly}).
\end{enumerate}
The trace class \(\mathcal{V}^{\mathrm{prim}}\) is the \(q_p \to 1\)
limit of \(\mathrm{SH}(q_p, t_p)\) in each prime factor; arithmetic
data sits in the tensor-indexing over primes, not in the deformation
parameters.
\end{theorem}
